% Generated draft: requires mathematical and source review.
\documentclass[11pt]{article}
\usepackage[margin=25mm]{geometry}
\usepackage{fontspec}
\setmainfont{Latin Modern Roman}
\usepackage{amsmath,amssymb,mathtools,unicode-math}
\setmathfont{Latin Modern Math}
\usepackage{enumitem,needspace,xurl,hyperref}
\setlength{\parindent}{0pt}
\setlength{\parskip}{5pt}
\setlength{\emergencystretch}{3em}
\newcommand{\N}{\mathbb{N}}
\newcommand{\Z}{\mathbb{Z}}
\newcommand{\Q}{\mathbb{Q}}
\newcommand{\R}{\mathbb{R}}
\newcommand{\C}{\mathbb{C}}
\newcommand{\F}{\mathbb{F}}
\newcommand{\A}{\mathbb{A}}
\newcommand{\PP}{\mathbb{P}}
\newcommand{\E}{\mathbb{E}}
\newcommand{\eps}{\varepsilon}
\newcommand{\dd}{\,\mathrm{d}}
\begin{document}
\section*{UnsolvedMath ID: 30002756}
\textbf{Source problem code:} catalogue entry
\subsection{Definitions and mathematical statement}
Under the usual Euclidean interpretation, let \(\Omega_1,\Omega_2\subset\R^d\) be Lipschitz domains in the same dimension. The supplied record does not specify the dimension, boundedness, or hypotheses ensuring discrete spectra. The formulation below is therefore conditional on the Laplacians having complete discrete eigenvalue sequences as presupposed by the source; boundedness is not added as a source hypothesis.

For \(i\in\{1,2\}\) and a constant \(\beta\geq0\), let \(A_{i,\beta}\) denote the Robin realization of \(-\Delta\) on \(L^2(\Omega_i)\), with boundary condition
\[
\partial_{\nu_i}u+\beta u=0\quad\text{on }\partial\Omega_i,
\]
where \(\nu_i\) is the outward normal and the condition is understood in the weak sense appropriate to a Lipschitz boundary. This uses the conventional nonnegative Robin parameter sign; \(\beta=0\) gives the Neumann condition. Let \(A_{i,D}\) denote the Dirichlet realization, with zero boundary trace. Write their eigenvalues in nondecreasing order, repeated according to multiplicity, as
\[
\bigl(\lambda_k(\Omega_i,\beta)\bigr)_{k\geq1},\qquad
\bigl(\lambda_k^D(\Omega_i)\bigr)_{k\geq1}.
\]
Here completeness of these sequences means that they exhaust the corresponding spectra.

The supplied question asks whether
\[
\begin{aligned}
\lambda_k(\Omega_1,\beta)&=\lambda_k(\Omega_2,\beta)
&&\text{for every }\beta\in[0,\infty)\text{ and every }k\geq1,\\
\lambda_k^D(\Omega_1)&=\lambda_k^D(\Omega_2)
&&\text{for every }k\geq1
\end{aligned}
\]
necessarily implies Euclidean congruence:
\[
\Omega_2=Q\Omega_1+a
\quad\text{for some }Q\in\R^{d\times d},\ Q^{\mathsf T}Q=I_d,
\text{ and }a\in\R^d.
\]
Equality is required at each common Robin parameter, together with Dirichlet spectral equality. The intended dimension and spectral hypotheses require source clarification before this conditional reconstruction can be regarded as a fully specified statement.
\subsection{Short English statement}
The question asks whether two Lipschitz domains with identical Robin spectra, including multiplicities, at every constant parameter \(\beta\geq0\), and identical Dirichlet spectra must be related by a Euclidean rigid motion. The supplied formulation leaves the dimension and assumptions guaranteeing complete discrete eigenvalue sequences unspecified. Its open status is supplied curation, not an independently verified conclusion.
\subsection{Sources}
\begin{itemize}
\item Local registry: \texttt{data/unsolvedmath/problems.json}, permanent UnsolvedMath ID \(30002756\); supplied source code \texttt{catalogue entry}, in the Oberwolfach Reports Open Problems collection. The supplied record gives the proposed year as 2015 and names no proposer.
\item Supplied report citation and literature source labeled ``original OWR report'': \emph{Spectral Theory and Weyl Functions} (2015), Oberwolfach Reports, pp.~5--86, DOI 10.4171/owr/2015/1. \url{https://doi.org/10.4171/owr/2015/1}. The supplied background describes this as a workshop report documenting open-problem discussions; no precise location within the report is supplied.
\item Supplied literature source labeled ``Arendt, ter Elst, and Kennedy, analytical aspects of isospectral drums'': \url{https://arxiv.org/abs/1305.1775}.
\item Literature assessment dated exactly \texttt{2026-08-21}: the supplied curation labels the problem open and states that the report explicitly poses the all-Robin-parameters plus Dirichlet spectral-family question as open. It further states that later work rules out a Robin intertwining for the standard counterexample but does not prove spectral-family congruence. This is supplied dated curation, not an independently verified literature assessment; the time-sensitive status and cited claims require independent verification.
\item The supplied statement status is \texttt{corrected\_verified}; its accompanying verification note says that mathematical content was preserved while notation, typography, or omitted definitions were repaired for publication. This is a registry assertion, not independent mathematical verification of this draft.
\end{itemize}
\end{document}
